\section{Introduction}\label{sec:intro}
In a ranked-choice election with $m$ candidates there are $m!$ rankings
(more with ties), so in an electorate of a few hundred a published ballot
is often unique and names its voter. A coercer exploits this in the
\emph{Italian attack}: they prescribe a distinctive ranking and check for
it on the board. Publishing the pairwise counts instead of the
ballots does not close the leak.  
In India, the votes of every polling station are counted
and published separately, so the result says how each village voted. In
the 2019 general election a Union minister told voters that villages
would be graded A to D by the share of their votes her party received,
and that development work would be allotted by grade~\cite{sabrang2019}.
The Law Commission of India had recommended in 2015 that the votes of
several stations be mixed in a totaliser before counting, so that no
station-level pattern is revealed~\cite{lawcom255}; the recommendation is
before the Supreme Court of India at the time of
writing~\cite{totaliser2026}. The finer the published result, the more
precisely a coercer can check how a group of voters behaved, and the
pairwise counts of a small ranked election are fine enough to check a
single voter. 
Every verifiable Schulze tally published so far reveals one or the other:
the shuffled ballots~\cite{haines2019verifiable} or the aggregate
matrix~\cite{vericondor2024}. The tally-hiding systems that reveal neither
use multi-party computation over threshold encryption, at a cost of hours
to days for twenty candidates~\cite{ordinosext2021,cgy2022}, or SNARKs
with a trusted setup~\cite{kryvos2022}.

We present a verifiable Schulze election that publishes neither the
ballots nor the pairwise counts and is built entirely from $\Sigma$-protocols, without MPC,
SNARKs, shuffles, or a trusted setup. A ballot is a matrix of encrypted
pairwise bits anchored to an encrypted position vector, so its
well-formedness is a conjunction of membership proofs on ciphertexts that
the verifier derives itself. The pairwise counts are aggregated homomorphically
and never decrypted on the board. The trustees certify the winner over
committed counts with a zero-knowledge certificate that the
strongest-path matrix is the closure of the defeat matrix. The
certificate has $O(m^3)$ proofs, a number that does not depend on the
electorate. In the terminology
of Kryvos~\cite{kryvos2022} the scheme is \emph{publicly} tally-hiding if
the trustees learn the pairwise counts. The public learns the winner and
the certificate's witness paths but no margin. 

Every proof the scheme uses comes from a certified library of
$\Sigma$-protocols in the Rocq theorem prover. The scheme's two
central properties are machine-checked on top of the library's theorems:
an accepted ballot encrypts a single ranking, and an accepted certificate
names a Schulze winner of the encrypted ballots. Because Rocq is a
constructive theorem prover, the client, the verifier, and the trustees'
prover are extracted from these proofs as executable code. The same
library, instantiated at Belenios's parameters, is a test oracle for
Belenios's verifier. It showed that an audit tool reported ``all checks
passed'' on an election containing an invalid ballot
(Section~\ref{sec:eval:oracle}).

\subsection{Contributions}
\begin{enumerate}
  \item \textbf{A publicly tally-hiding Schulze election from
  $\Sigma$-protocols alone} (Section~\ref{sec:schulze}). Ballots are
  encrypted pairwise bits anchored to an encrypted position vector, so
  well-formedness is a conjunction of membership proofs and needs no
  shuffle. Pairwise counts are aggregated homomorphically and never decrypted on
  the board, and the winner is certified over committed counts by a
  zero-knowledge closure certificate of $O(m^3)$ proofs, a number that
  does not depend on the electorate. It needs no trusted setup,
  where Kryvos needs Groth16. It certifies a $10$-candidate, $50$-voter
  election in $18$ minutes (verification $9$ minutes) and a
  $17$-candidate, $96$-voter one in $2.4$ hours (verification $71$
  minutes). MPC-based tally-hiding Schulze takes hours for twenty
  candidates and $1024$ voters~\cite{cgy2022} or days for twenty
  candidates~\cite{ordinosext2021}.
  \item \textbf{Machine-checked ballot and certificate soundness}
  (Section~\ref{sec:mech}). \rocq{ballot_sound}: an accepted ballot
  encrypts a single ranking. \rocq{verify_tally_sound}: an accepted
  certificate names a Schulze winner of the encrypted pairwise counts.
  \rocq{ballot_zk} and \rocq{tally_zk}: every ballot and certificate
  transcript is perfectly simulatable from the public ciphertexts and,
  for the certificate, the witness paths. \rocq{schulze_end_to_end}:
  accepted ballots and an accepted certificate together name a Schulze
  winner of the rankings the voters cast. 
  The extracted client and verifier from Rocq proofs run on the $13$ Gentoo
  council elections of 2011--2023: every published ballot is accepted,
  every certificate verifies, and every certified winner is the
  first-seat winner of Gentoo's official count
  (Section~\ref{sec:eval:gentoo}).
  \item \textbf{A certified $\Sigma$-protocol library}
  (Section~\ref{schn_proof}): completeness, special soundness, and special
  honest-verifier zero-knowledge for the Schnorr, Chaum--Pedersen,
  encryption, and decryption proofs and their $n$-ary compositions,
  generic in the group and extracted to executable code. Used as a test
  oracle for Belenios's verifier, the extracted code found a defect in an
  audit tool (Section~\ref{sec:eval:oracle}). The Rocq proofs are at
  \textcolor{rocqlink}{\url{https://anonymous.4open.science/r/SigmaProtocol-EBA5/README.md}},
  and their documentation is at
  \textcolor{rocqlink}{\url{https://anonymous.4open.science/r/SigmaProtocol-EBA5/doc/html/toc.html}}.
\end{enumerate}

Section~\ref{sec:background} gives background and
Section~\ref{schn_proof} describes the library. Then come the scheme
(Section~\ref{sec:schulze}), its mechanisation (Section~\ref{sec:mech}),
the evaluation (Section~\ref{sec:eval}), a discussion of what is not
covered (Section~\ref{sec:discussion}), and related work. Rocq names and terms
in \textcolor{rocqlink}{blue} link to their definitions in the artefact.
